\section{مقدمه}
\label{sec:intro}

\subsection{مسئله و انگیزه}

مجموعه‌داده‌ی \lr{IMDb} \cite{imdb, imdbdocs} یکی از بزرگ‌ترین مجموعه‌داده‌های رایگان در زمینه‌ی فیلم و سریال است. هفت جدول دارد که در مجموع نزدیک به \lr{190} میلیون سطر می‌شود (فقط جدول \lr{title\_principals} حدود \lr{90} میلیون سطر دارد). هدف این پروژه این است که یک پایگاه‌داده‌ی \lr{PostgreSQL 16} بسازیم که بتواند این حجم داده را هم از نظر فضای دیسک و هم از نظر سرعت کوئری، خوب مدیریت کند.

اگر بخواهیم فایل‌های \lr{TSV} فشرده را مستقیم و ساده بارگذاری کنیم (مثلاً با \lr{INSERT} تک‌تک یا حتی \lr{COPY} ساده)، دو مشکل پیش می‌آید:

\begin{enumerate}
\item \textbf{فضای دیسک.} اگر هر پیام فقط یک سطر باشد، اندازه‌ی آن از آستانه‌ی \lr{TOAST} (حدود \lr{2} کیلوبایت) رد نمی‌شود، یعنی ستون‌های \lr{JSONB} بدون فشرده‌سازی ذخیره می‌شوند. برای \lr{190} میلیون سطر، این چند ده گیگابایت فضای اضافه می‌شود که با محدودیت دیسک ما (حدود \lr{52} گیگابایت) جور در نمی‌آید.
\item \textbf{هم‌زمانی نوشتن.} می‌خواهیم چند پردازش هم‌زمان داده بنویسند تا سریع‌تر شود، اما اگر از قفل معمولی (\lr{FOR UPDATE} بدون \lr{SKIP LOCKED}) استفاده کنیم، همه پشت سر هم منتظر می‌مانند.
\end{enumerate}

راه‌حلی که انتخاب کردیم، گذاشتن یک صف پیام بین دانلود داده و درج نهایی است. یک تولیدکننده (\lr{producer}) فایل‌های فشرده را جریانی می‌خواند (بدون این‌که کامل روی دیسک استخراج شود)، سطرها را در دسته‌های \lr{1000}تایی بسته‌بندی و در صف می‌گذارد. یک یا چند مصرف‌کننده (\lr{consumer}) هم با \lr{SELECT ... FOR UPDATE SKIP LOCKED} \cite{pglocking} دسته‌ها را بدون رقابت با هم برمی‌دارند و با \lr{execute\_values} به‌صورت انبوه درج می‌کنند.

\subsection{اهداف گزارش}

این گزارش سه هدف دارد:

\begin{itemize}
\item توضیح کامل معماری پروژه: از تنظیمات \lr{PostgreSQL} و طراحی جدول‌ها تا صف پیام و خط لوله‌ی بارگذاری؛
\item نمایش و بررسی هشت سناریوی کوئری که سند پروژه خواسته، هم قبل از ایندکس‌گذاری و هم بعد از آن؛
\item پاسخ به سؤال‌های نظری پروژه درباره‌ی تنظیمات \lr{PostgreSQL}، انواع ایندکس، \lr{JSONB}/\lr{TOAST} و نظریه‌ی صف پیام.
\end{itemize}

\subsection{ساختار گزارش}

بخش \ref{sec:architecture} معماری اجرا و طراحی اسکیما را توضیح می‌دهد. بخش \ref{sec:queue} صف پیام و کد تولیدکننده/مصرف‌کننده را نشان می‌دهد. بخش \ref{sec:scenarios} هشت سناریوی کوئری را معرفی می‌کند. بخش \ref{sec:indexing} ایندکس‌ها و روش مقایسه‌ی عملکرد را توضیح می‌دهد. بخش \ref{sec:theory} به سؤال‌های نظری پاسخ می‌دهد و بخش \ref{sec:conclusion} جمع‌بندی است.
